\documentclass[11pt,a4paper]{article}
\usepackage[margin=22mm]{geometry}
\usepackage{amsmath,amssymb}
\setlength{\parskip}{5pt}
\setlength{\parindent}{0pt}
\widowpenalty=10000
\clubpenalty=10000
\brokenpenalty=10000
\pagestyle{myheadings}
\markright{Click2Future, Group 7. Validation pilots; final experiments pending.}
\title{Click2Future: Action-Conditioned Video Generation for Tabletop Futures}
\author{Group 7, AI6132 Generative AI For Visual Synthesis\\
\small PRAMOD SANJAY, KAPOOR ANANYA,\\
\small MAITI SHRIJEET, BOBAN ATHUL KRISHNA}
\date{\small Working draft; measured validation pilots}
\begin{document}
\maketitle
\begin{abstract}
We study whether a compact generative video model can predict several physical futures while following an explicit object selection and force. A procedural tabletop simulator provides initial images, action inputs, trajectories and contact labels. The implemented methods comprise a conditional video variational autoencoder and a small pixel diffusion model with numeric or spatial action conditioning. Evaluation separates image appearance, tracked motion, collisions and ensemble uncertainty. We train E0, E2, E3 and E4 on 336 synthetic clips and evaluate 24 validation conditions from four held-out scenes. None produces complete tracked trajectories. A flow-decoder CVAE improves object visibility but remains stationary under force controls. A fixed-checkpoint 16/32/64-step DDIM study also fails the video-quality gate. These negative results expose the limits of the current models and appearance metrics; they do not establish a successful learned physical model or a final test-set ranking.
\end{abstract}

\section{Introduction}
An image can show the location and shape of an object without revealing its mass or friction. Applying the same visible force can therefore produce several different motions. A useful video generator should represent this uncertainty while preserving a clear dependence on the requested action. Four visually different samples are insufficient evidence: the selected object must respond, unselected objects must behave consistently with contacts, and differences between futures should resemble variation permitted by the simulator.

Our task is to predict a distribution over future RGB frames given an initial image, a selected-object mask and a force vector. The environment is a bounded, top-down tabletop with rigid objects and simple appearance. This restriction makes physical labels available and keeps training feasible on the course-provided TC2 cluster. We do not study robot control, arbitrary photographs or general three-dimensional dynamics.

The main comparison is a conditional video VAE (E0) against two versions of video diffusion: numeric conditioning (E2) and spatial conditioning (E3). First-frame-only diffusion (E1) is an action-blind control. Auxiliary trajectory and collision supervision (E4) is an optional ablation after the base generator works. The review question is whether localizing the requested action improves generated physical behavior under comparable training settings. No novelty or improvement is inferred from implementing this design.

Variational autoencoders use a learned approximate posterior to optimize latent-variable models~\cite{kingma}. Conditional variants model structured outputs given an observed input~\cite{sohn}. Denoising diffusion models instead learn to reverse a gradual corruption process~\cite{ho}. DDIM provides a sampling procedure that permits a reduced number of denoising steps~\cite{song}. We implement compact versions of these approaches for our own tabletop task. Published image-generation results are not used as evidence for this project's video performance.

\section{Methods}
\subsection{World, action and observations}
The simulator uses a unit square with image coordinates: positive $x$ points right and positive $y$ points down. A constant force acts at the selected object's centre for 0.25 seconds, followed by unforced motion. It is a force applied over a duration, rather than an instantaneous velocity change. Clicking elsewhere on an object does not introduce an off-centre torque. Wall contacts and object contacts follow the simulator's rigid-body implementation. Friction and restitution act at contacts; there is no separate tabletop drag. Walls retain frictionless, zero-elasticity defaults.

Core clips contain 32 frames at 8 frames per second and $64\times64$ resolution. Frame zero precedes the action; the 31 predicted future frames extend to 3.875 seconds. Physics advances at 240 Hz. Each record contains RGB frames, object masks, positions, valid object slots, action metadata and contact onsets aggregated over the intervals between rendered frames. Replay uses the saved scene specification and simulator settings. The 128px extension is deferred until coherent 64px generation and the main comparisons are established.

Mass, contact friction and restitution are sampled independently of visible appearance. Multiple references for the same observable condition share frame zero, the selected mask, force and timing while varying hidden physics. All actions and physics variants of a visible scene remain in one split. This grouping prevents a nearly identical initial scene from appearing in both training and evaluation.

The controlled corpus includes zero force, a force and its reverse, two collision directions and a wall interaction. A separate generalization corpus defines disjoint training, validation, ID-test, shape-OOD, texture-OOD, force-OOD and repeated-physics splits. Its triangles and object-attached patterns differ from the earlier corpus. Final ID/OOD claims require models trained on that same corpus; testing an earlier model on a changed corpus does not isolate one factor.

Only frame-zero RGB, the selected mask and numeric force enter the model. Oracle future masks, positions, contacts and hidden physical coefficients are unavailable at inference. RGB is normalized to $[-1,1]$. Numeric forces are divided by 0.5 without clipping; a separately declared physical magnitude limit controls accepted inputs. The spatial action field is
\[
 A_c(u,v)=M(u,v)f_c,\qquad c\in\{x,y\},
\]
where $M$ is the selected mask and $f$ is the action vector. Using the same selected mask in E2 and E3 prevents object identification alone from explaining their difference.

\subsection{Conditional video VAE}
E0 has a learned conditional Gaussian prior $p_\psi(z\mid c)$ and posterior $q_\phi(z\mid c,Y)$, where $c$ denotes the observable condition and $Y$ the future clip. A 2D convolutional encoder processes the initial image, mask and broadcast force. A 3D encoder processes future RGB only during posterior training. The decoder receives condition features, a global clip latent and explicit normalized time fields. Its output is a complete future clip, bounded by a tanh activation.

The implemented pilot has 574,659 parameters, base width 16 and latent dimension 64. Its loss combines future-only mean absolute reconstruction error with mean per-latent Gaussian KL divergence:
\[
 \mathcal L_{\mathrm{VAE}}=\mathbb E_{z\sim q_\phi}\!\left[\|\hat Y-Y\|_1\right]
 +\beta\,\mathbb E_{j}\!\left[D_{\mathrm{KL}}(q_{\phi,j}\|p_{\psi,j})\right].
\]
The beta warmup depends on the saved absolute training step. Sampling uses the conditional prior, with one private random generator per seed. The posterior never runs during prior inference. Validation posterior reconstruction and actual prior-sample quality are reported separately: a good posterior score need not imply a useful sampled future.

\subsection{Pixel video diffusion}
The implemented diffusion model operates directly on RGB future frames. It is not a latent video diffusion system; the proposed frame autoencoder is outside the current critical path. The denoiser combines spatial convolutions, temporal convolutions, a spatial encoder-decoder with skip connections, and temporal attention at reduced spatial resolution. The initial frame and action channels are broadcast over future time. Diffusion timesteps and numeric actions provide conditioning embeddings.

For a clean future $Y$ and Gaussian noise $\epsilon$, the forward process is
\[
 Y_t=\sqrt{\bar\alpha_t}Y+\sqrt{1-\bar\alpha_t}\epsilon.
\]
The code supports either epsilon prediction or velocity prediction,
$v_t=\sqrt{\bar\alpha_t}\epsilon-\sqrt{1-\bar\alpha_t}Y$, with mean squared target error. Main ablations must use the same selected target, schedule, backbone width and initialization policy. Comparing a velocity model to an epsilon model would confound the conditioning comparison.

E1 replaces the selected mask and action with zero inputs. E2 supplies the selected mask and broadcasts numeric action fields over the image. E3 multiplies those fields by the selected mask, retaining the same backbone dimensions and numeric embedding. This is the implemented object-centric conditioning mechanism. E3 does not add a learned object detector or receive oracle future object states.

Inference uses EMA weights and deterministic DDIM updates with private seed generators. Clean estimates are clipped to $[-1,1]$ and epsilon is recomputed after clipping. The exact original first frame is prepended to every generated future. A 16/32/64-step validation study is reported below; no setting passes the physical-video checks.

\subsection{Auxiliary supervision}
E4 adds trajectory and contact heads to denoiser features. Training targets use a canonical frame-zero object order and masks for valid objects and pairs. Position estimates use spatial heatmap moments; contact outputs are pair logits. Trajectory and contact losses are applied according to the configured timestep gate and weights. Those settings remain starting values until validated.

The head outputs are additional predictions, not inference conditioning. Auxiliary trajectories can look accurate even when generated RGB is poor. Consequently, head scores are exported separately from trajectories and contacts extracted from the generated video. A comparison of E3 and E4 must assess decoded-video behavior as well as the supervised heads.

\section{Experiments and analysis}
\subsection{Frozen evaluation and failure handling}
The evaluator freezes complete held-out condition lists, every reference future, ordered sample seeds and metric settings. Dataset, split, reference, configuration and evaluator-code hashes bind a plan to its inputs. The checkpoint runner checks method identity and training-corpus compatibility, produces four futures for every case and binds checkpoint/video hashes to the prediction manifest. Partial runs cannot produce a successful final summary.

A palette-based tracker uses object identities from the initial frame. It is designed for the synthetic appearance domain and can fail on blur, missing objects, occlusion or unfamiliar texture. Missing RGB and tracking failures remain in the denominators. Conditional averages and undefined rates are accompanied by counts; a completely missing prediction does not receive zero error or a favourable action score.

Trajectory evaluation reports ADE and FDE, both for all objects and for the selected object. Mean-over-samples and whole-clip best-of-four scores remain separate. Collision evaluation distinguishes any encounter during the clip, exact interval labels and one-to-one timed onset matches. High encounter accuracy does not establish correct timing. SSIM and LPIPS measure appearance against one declared reference, while trajectory coverage and ensemble spread use the repeated-physics references. Diversity is assessed alongside quality and adherence, rather than treated as an independent success measure.

Aggregation first averages dependent samples and conditions within a visible scene and training run, then averages training runs and independent scenes. Confidence intervals resample whole scene groups. A single scene group has no bootstrap interval. The planned independent training repeats and final dataset size have not yet been completed or frozen.

\subsection{Action controls and extractor calibration}
Zero-force and reversed-force checks use the first 0.5 seconds, ending at frame four, with a fixed one-pixel motion threshold at 64px. Every hidden-physics reference must start at rest and avoid contact during this window. Eligibility is decided from references before model output is inspected. Excluded groups remain listed with their reasons.

Selected-object controls require a separate intervention: the initial image, physical objects, force vector and duration remain identical under two different clicks. Derived references live outside the training corpus. Correct target following requires selected motion along the force and no above-threshold excursion by unselected objects in both conditions. A separate per-object contrast rejects identical motion under both clicks. The E1 control produces identical videos for both selections when seeds match, as expected for an action-blind model.

In a two-scene validation correctness check, the original 0.15 force left all eight exact-label target pairs below the one-pixel criterion, whereas rasterized references crossed it more often. This exposes quantization sensitivity rather than a model failure. That calibration was retained. A separate shared-magnitude 0.4 probe kept the threshold unchanged and passed for both continuous-label and rendered-reference controls. Static and ignored-click controls failed. These values describe a simulator validation check, not the performance of any trained method. A magnitude used for final model tests must be frozen using validation beforehand.

Contact extraction has an independent timing limitation. On a fresh 63-clip controlled check over three scene groups, all 3,906 object observations were tracked. Any-encounter scoring produced TP=23, FP=0, FN=1 and TN=39, while timed-event F1 was 0.551. The stronger occurrence result cannot substitute for the weaker event-timing result. These extractor measurements do not count the 63 dependent clips as 63 independent scenes.

\subsection{Multi-scene validation pilot}
Dataset-0.3 seed 730000 contains 16 training scenes (336 clips), four validation scenes (84 reference clips) and 248 reserved test clips across the ID/OOD/physics splits. All 668 clips replayed exactly before training. Validation has six observable action conditions per scene, including zero, forward, reverse, both collision directions and a wall interaction. Four sample seeds per condition give 96 videos per method. Repeated hidden-physics references support coverage analysis; image distances use one explicitly declared reference per condition.

TC2 job 43291 independently trained the RGB CVAE for 20,000 updates (batch two, seed zero), and E2/E3/E4 for 5,000 updates each (batch one, seed seven). All diffusion methods use width 32, FP32 velocity targets, EMA decay 0.995, 1,000 noise steps and 32 DDIM sampling steps. The CVAE uses width 16, latent size 64 and beta 0.001 with 20-percent warmup. Both losses use foreground weight 32 with a two-pixel margin. Future masks weight training errors only. Validation posterior loss selected the CVAE checkpoint at update 16,000; diffusion uses its final EMA checkpoint. The initial four-method training allocation lasted 2,087 seconds; its evaluation allocation lasted 372 seconds on an A40.

The common diffusion initialization policy uses independent runs with no inherited training. E2 and E3 each have 964,035 parameters; E4 adds 2,580 head parameters. Extra E4 initialization consumes random draws, so the same seed does not assert identical E3/E4 training noise. Training-only contact counts (156 positive and 10,260 negative unordered pair intervals) set the capped positive weight to 20. One training seed, unequal CVAE/diffusion budgets and only four validation scenes prevent a final method ranking.

\begin{table}[!ht]
\centering
\small
\begin{tabular}{lrrrr}
\hline
Method & SSIM & LPIPS & Tracked / 5,952 & Complete / 96 \\
\hline
E0 RGB CVAE & 0.9291 & 0.1230 & 398 & 0 \\
E2 numeric diffusion & 0.4736 & 0.1503 & 378 & 0 \\
E3 spatial diffusion & 0.4806 & 0.1493 & 269 & 0 \\
E4 auxiliary diffusion & 0.4711 & 0.1537 & 162 & 0 \\
Static frame-zero copy & 0.9446 & 0.0622 & n/a & n/a \\
\hline
\end{tabular}
\caption{Validation pilot on 24 conditions and four scenes. SSIM is higher-is-better; LPIPS is lower-is-better. Image scores exclude frame zero and average cases within scene, then scenes. Tracking uses candidate extractor version 2; complete means all valid future object tracks. Static copying is an appearance control, not a successful future predictor.}
\label{tab:multiscene}
\end{table}

No method has a complete selected-object sample, so mean and best-of-four ADE/FDE, trajectory spread and reference coverage are undefined. Undefined values remain missing rather than becoming zero error. All 16 forward, reverse and zero-action samples per method lack the required selected tracks; all 16 selected-object intervention pairs also fail tracking. These outcomes do not support an action-adherence or meaningful-uncertainty claim.

The whole-scene 95-percent LPIPS bootstrap intervals are [0.0797, 0.1746] for E0, [0.1404, 0.1601] for E2, [0.1394, 0.1579] for E3 and [0.1425, 0.1649] for E4. They are conditional on the supplied checkpoints and do not include training-seed uncertainty. The static-copy LPIPS interval is [0.0460, 0.0783]. Its stronger appearance scores expose background and displacement sensitivity: good single-reference image distance does not establish useful physical prediction.

The RGB CVAE's validation posterior MAE is 0.02091 and prior MAE 0.02113 in normalized RGB units. Training posterior/prior errors are 0.01205/0.01386. Its validation mean per-latent KL is 0.2741 and pairwise prior RGB MAE is 0.00330. Nonzero KL and sample variation do not establish useful latent uncertainty; these diagnostics alone also do not justify calling the failure posterior collapse.

\subsection{Auxiliary heads versus generated RGB}
The E4 analysis replays all 96 validation samplings while capturing the actual final timestep-zero DDIM heads. Every RGB pixel matches the saved sample. Head predictions have all-object ADE 0.1314 and FDE 0.1560 in normalized coordinates against one declared reference. At the fixed 0.5 probability threshold, timed-contact matching with one-interval tolerance gives TP=28, FP=331 and FN=36, hence event F1=0.1324. Exact-interval positive prevalence is 64/2,976 and interval F1=0.1087. Any-encounter F1=0.5714 does not establish correct timing.

Only 162 of 5,952 future object observations are jointly available from the head and RGB extraction. Their mean position disagreement is 0.1481. No pair interval is jointly observed, so RGB/head collision agreement is undefined. Accurate-looking auxiliary outputs cannot substitute for decoded-video behavior. A separate 5,000-update check on the first sorted training collision clip has head ADE 0.00493 and contact-event F1 1.0 (TP=4, FP=0, FN=0 across four samples). RGB tracks 247/248 observations, with three complete clips. One missing selected observation still fails the strict gate; conditional selected ADE 0.00396 over the other three clips is not a best-of-four score. This training-only result shows that the head can learn a known contact, not that physical behavior generalizes. The allocation lasted 423 seconds.

\subsection{Validation sampling-step study}
A predeclared study holds the E3 checkpoint, sample seeds and metric settings fixed. It uses the first four sorted validation scenes and zero, forward, reverse and collision-0-to-1 conditions, for 16 conditions and 64 videos per setting. Excluded validation conditions remain listed. This is a subset study rather than the full 24-condition comparison.

\begin{table}[!ht]
\centering
\begin{tabular}{rrrr}
\hline
DDIM steps & Export seconds & Tracked / 3,968 & SSIM \\
\hline
16 & 16.25 & 330 & 0.5473 \\
32 & 29.37 & 220 & 0.5083 \\
64 & 57.59 & 184 & 0.4931 \\
\hline
\end{tabular}
\caption{Same frozen E3 checkpoint on the fixed validation subset. Export time includes checkpoint loading and writing all 64 videos on a TC2 A40. It is not a pure denoiser latency measurement.}
\label{tab:steps}
\end{table}

All three settings have zero complete selected-object clips and fail the action checks through missing tracks. More sampling steps increase cost without repairing this checkpoint. The study does not select a successful sampler or establish a general advantage for fewer steps. The allocation lasted 160 seconds, within the reserved 0.5-hour study limit.

\subsection{CVAE decoder follow-up}
The RGB samples change object appearance on held-out scenes. A separate exploratory CVAE therefore predicts a backward displacement field and bilinearly samples observed frame-zero pixels. Its zero-initialized output copies the initial image. The latent distributions, data, seed, RGB/KL loss and 20,000-update limit are retained, but decoder parameterization and initialization change. The model has 574,226 parameters. It cannot reconstruct unseen textures, enforce rigidity or enforce collision laws.

This flow model tracks 5,654/5,952 observations (94.99 percent), with 80/96 complete clips. Its single-reference SSIM is 0.9693 and LPIPS 0.0227. Nevertheless all 16 forward samples, all 16 reverse samples and all 16 selected-object intervention pairs are stationary in the fixed action window. Correct force-sign and target-following rates are zero over complete eligible pairs. Native-force interpretation is limited: 11/16 continuous and 8/16 rendered reference pairs also fall below the short-window motion threshold. In contrast, the separate force-0.4 selection probe passes all 16 continuous and rendered reference pairs, while this CVAE passes none. Better appearance is therefore insufficient. A further displacement-supervised variant is kept separate from the ordinary CVAE: future masks and positions enter a training loss only, not generator conditioning. The displacement-supervised run completes 20,000 updates, with its combined validation criterion selecting update 9,000. It tracks 5,765/5,952 observations and 80/96 complete clips, with SSIM 0.9729 and LPIPS 0.0168. Yet only one of 16 reversal pairs has the correct force signs: 15 forward samples move, while only one reversed sample moves. All 16 selected-object pairs remain stationary. Additional supervision improves these appearance scores but does not establish reliable action control. Neither exploratory variant provides an independent final test.

\clearpage
\begin{figure}[!ht]
\centering
\pdfximage width\linewidth{report/figures/multiscene_validation.png}\pdfrefximage\pdflastximage
\caption{First sorted validation forward condition, seed zero, at six fixed times. The simulator is one declared reference. E0 uses its conditional prior; E2/E3/E4 use EMA and 32 DDIM steps. Both flow variants preserve appearance but fail the selection controls. The supervised variant also fails 15/16 force-reversal pairs. Selection is fixed before visual inspection; this is not a best-sample comparison.}
\label{fig:multiscene}
\end{figure}
\clearpage
\section{Conclusion}
The simulator and evaluation workflow support a controlled study of action-conditioned tabletop futures. The implemented baselines make the intended numeric-versus-spatial comparison possible, and reference calibration exposes failure modes that simple image diversity or collision accuracy can hide. The four-method validation pilot and sampling-step study retain failed tracking and action response. The flow follow-up shows that appearance can improve without force adherence. These results do not establish whether spatial conditioning or auxiliary supervision improves useful generated futures. Improving coherent, action-responsive validation generation is necessary before opening the reserved test suites or claiming physical uncertainty. Possible extensions include a validated latent representation and 128px generation after the core 64px comparison works.

\begin{thebibliography}{9}
\bibitem{kingma} Diederik P. Kingma and Max Welling. Auto-Encoding Variational Bayes. arXiv:1312.6114, 2013.
\bibitem{sohn} Kihyuk Sohn, Honglak Lee and Xinchen Yan. Learning Structured Output Representation using Deep Conditional Generative Models. Advances in Neural Information Processing Systems 28, 2015.
\bibitem{ho} Jonathan Ho, Ajay Jain and Pieter Abbeel. Denoising Diffusion Probabilistic Models. arXiv:2006.11239, 2020.
\bibitem{song} Jiaming Song, Chenlin Meng and Stefano Ermon. Denoising Diffusion Implicit Models. ICLR, 2021.
\end{thebibliography}
\end{document}
